%%% -*- coding: utf-8 -*-
% !TeX root = main-slides.tex

\lecture{Reflection and Annotations}{reflection}
\section{Reflection and Annotations}

\showtitlepage{Reflection and Annotations}

%%% Every example is taken from SEE, so the sources named in the listings can
%%% be opened next to the slide.

\begin{frame}{Why Reflection at All?}
  Normally the compiler resolves every name for you: which class, which field,
  which method, all fixed when the code is built. Sometimes it cannot. A
  message arrives over the network and names its own class. A configuration
  file names a tool. A menu is built from whatever a city happens to offer, and
  a test runner must find methods that were written long after it was.
  \vspace{1ex}

  \begin{itemize}
    \item \emph{Reflection} is the run-time view of that metadata: ask an
      object what it is, list its members, create instances, call methods ---
      all decided while the program runs.
    \item \emph{Annotations}, in C\# called \emph{attributes}, are the notes
      you leave in that metadata for such a reader to find.
    \item Both live at the edges of a program: frameworks, factories, editors,
      serialisation. In the middle of your code they are the wrong tool.
  \end{itemize}
  \vspace{2ex}

  \hint{SEE sends an action over the network as its type name plus JSON, and
  the receiver builds the object from that name --- there is no other way to
  know which action arrived.}
\end{frame}

\begin{frame}[fragile]{The \texttt{Type} Object}
  \framelabel{slide:typeobject}
  \begin{lstlisting}[basicstyle=\scriptsize\ttfamily]
SEECity city = FindFirstObjectByType<SEECity>(); // from the scene, not new
Type declared = typeof(SEECity);        // known when the code is built
Type actual   = city.GetType();         // whatever the object really is
Type named    = Type.GetType(typeName); // from a string; null if unknown

FieldInfo[] fields = actual.GetFields();         // public instance ones
MethodInfo  reset  = actual.GetMethod("Reset");  // public, by name

FieldInfo hidden = actual.GetField("graph",            // everything else has
    BindingFlags.NonPublic | BindingFlags.Instance);   // to be asked for
  \end{lstlisting}
  \vspace{1ex}

  \begin{itemize}
    \item \texttt{Type} is the entry point, and every member comes back as an
      object of its own --- \texttt{FieldInfo}, \texttt{MethodInfo},
      \texttt{PropertyInfo} --- carrying its name, its type and its
      annotations.
    \item \texttt{BindingFlags} decides what you are shown. The default is
      public and instance; asking for anything else and forgetting to say so
      is not an error but a \texttt{null}.
  \end{itemize}
\end{frame}

\begin{frame}[fragile]{From a Type to an Object}
  \begin{lstlisting}[basicstyle=\scriptsize\ttfamily]
// Assets/SEE/DataModel/DG/Attributable.cs -- clone without knowing what
Attributable clone = (Attributable)Activator.CreateInstance(GetType());

// Assets/SEE/Net/Actions/AbstractNetAction.cs -- the type travels as text
string data = action.GetType().ToString() + ';' + JsonUtility.ToJson(action);
...
string[] tokens = data.Split(new[] { ';' }, 2, StringSplitOptions.None);
AbstractNetAction received = (AbstractNetAction)JsonUtility.FromJson(
    tokens[1], Type.GetType(tokens[0]));
  \end{lstlisting}
  \vspace{1ex}

  \begin{itemize}
    \item \texttt{Activator.CreateInstance} calls the parameterless
      constructor: there has to be one, and it has to be accessible.
    \item \texttt{CreateInstance} and \texttt{FromJson} both hand back an
      \texttt{object}: the cast in front --- to \texttt{Attributable}, to
      \texttt{AbstractNetAction} --- is the only check you get, and it happens
      when the line runs.
  \end{itemize}
\end{frame}

\begin{frame}[fragile]{Calling a Method by Name}
  \begin{lstlisting}[basicstyle=\scriptsize\ttfamily]
// Assets/SEE/Game/City/AbstractSEECity.cs -- all four are declared there
reset.Invoke(city, null);                       // Reset() is void: null

MethodInfo save = actual.GetMethod("Save");     // void Save(string filename)
save.Invoke(city, new object[] { "city.cfg" }); // arguments in one array

MethodInfo types = actual.GetMethod("HierarchicalEdgeTypes");
object all = types.Invoke(null, null);          // static: no receiver
// two overloads are public, so GetMethod must be told which one
MethodInfo hover = actual.GetMethod("UserIsHoveringCity", Type.EmptyTypes);
bool yes = (bool)hover.Invoke(city, null);      // a boxed bool
  \end{lstlisting}
  \vspace{1ex}

  \begin{itemize}
    \item The receiver comes first --- \texttt{null} for a static method,
      otherwise an instance of the declaring type, or a
      \texttt{TargetException}.
    \item The arguments come as one \texttt{object[]}, not as a parameter list.
      \texttt{null} means none, optional parameters are not filled in for you,
      and a wrong count is a \texttt{TargetParameterCountException}.
    \item The result is an \texttt{object}: \texttt{null} for \texttt{void}, a
      boxed copy for a value type, so the cast must name the exact type ---
      \texttt{(long)} on a boxed \texttt{int} throws.
  \end{itemize}
\end{frame}

\begin{frame}[fragile]{Annotations Are Data}
  \framelabel{slide:attributes}
  \begin{lstlisting}[basicstyle=\scriptsize\ttfamily]
// Assets/SEE/Utils/EnvironmentVariable.cs
[AttributeUsage(AttributeTargets.Field)]      // where it may be written
public class EnvironmentVariableAttribute : Attribute
{
  public string VariableName { get; }
  public EnvironmentVariableAttribute(string variableName) { ... }
}

[EnvironmentVariable("DEBUG_LOG")]            // a constructor call, stored
public bool debugLog;                         // in the metadata
  \end{lstlisting}
  \vspace{1ex}

  \begin{itemize}
    \item An attribute is an ordinary class derived from \texttt{Attribute}.
      Writing it in brackets records a call to its constructor; it runs when
      somebody asks for the attribute, never when the field is used.
    \item The suffix is convention: the class is \texttt{XAttribute}, the
      annotation is written \texttt{[X]}.
    \item You have seen them all along: \texttt{[SerializeField]},
      \texttt{[Header]} (Unity), \texttt{[Test]} (NUnit),
      \texttt{[ShowInInspector]} (Odin), \texttt{[Obsolete]} (the compiler
      itself reads that one).
  \end{itemize}
\end{frame}

\begin{frame}[fragile]{Reading Annotations}
  \begin{lstlisting}[basicstyle=\scriptsize\ttfamily]
// one of them, if it is there at all -- inherit is true unless you say so
RuntimeTabAttribute tab = member.GetCustomAttribute<RuntimeTabAttribute>();
RuntimeTabAttribute own =
    member.GetCustomAttribute<RuntimeTabAttribute>(inherit: false);

// all of them, to filter -- Assets/SEE/UI/RuntimeConfigMenu/RuntimeTabMenu.cs
if (member.GetCustomAttributes().Any(a => a is ObsoleteAttribute)) { ... }
string name = member.GetCustomAttributes().OfType<RuntimeTabAttribute>()
                    .FirstOrDefault()?.Name;
  \end{lstlisting}
  \vspace{1ex}

  \begin{itemize}
    \item The attribute object is constructed when you ask for it, from the
      call recorded in the metadata. Asking twice gives you two objects.
    \item The \texttt{inherit} parameter decides whether an annotation on a
      base member counts as one on the member you are looking
      at.\footnote{\texttt{MemberInfo.GetCustomAttributes(bool)} ignores it for
      properties and events; the \texttt{Attribute} helpers behind the
      extension methods above do emulate it.}
    \item Nothing at all happens without a reader. Unity's inspector, NUnit's
      runner and SEE's runtime configuration menu are such readers --- and in
      the next exercises, so are you.
  \end{itemize}
\end{frame}

\begin{frame}[fragile]{What Odin Is}
  \begin{lstlisting}[basicstyle=\scriptsize\ttfamily]
// Assets/SEE/Game/City/AbstractSEECity.cs -- Odin's base class, not Unity's
public abstract partial class AbstractSEECity : SerializedMonoBehaviour

// Assets/SEE/Game/City/SEECity.cs -- what Unity's serializer cannot store
[OdinSerialize, ShowInInspector, TabGroup(DataFoldoutGroup), ...]
public SingleGraphPipelineProvider DataProvider = new();
  \end{lstlisting}
  \vspace{1ex}

  \begin{itemize}
    \item Odin Inspector (by Sirenix) is a bought editor
      extension, in \texttt{Assets/Plugins/Sirenix} like every third-party
      asset. It replaces Unity's inspector and brings a serializer of its own.
    \item SEE uses both halves: a city is a \texttt{SerializedMonoBehaviour},
      so dictionaries, interfaces and polymorphic objects survive a save; and
      its hundred-odd settings are laid out by annotations ---
      \texttt{[TabGroup]}, \texttt{[Button]}, \texttt{[ShowIf]}.
    \item Only half of that ships:
      \texttt{Sirenix.OdinInspector.Attributes.dll} is a run-time assembly and
      is in the player, while \texttt{Sirenix.OdinInspector.Editor.dll}, which
      reads those annotations and draws them, is not.
  \end{itemize}
\end{frame}

\begin{frame}[fragile]{Odin's Annotations at Run Time}
  Odin's annotations are still in the running game --- their assembly ships
  with it. The code that reads them is not, so SEE writes its own reader, over
  twin annotations of its own:
  \vspace{1ex}

  \begin{lstlisting}[basicstyle=\scriptsize\ttfamily]
// Assets/SEE/Game/City/SEECity.cs -- Odin on the left, our twin on the right
[TabGroup(DataFoldoutGroup), RuntimeTab(DataFoldoutGroup), ShowInInspector]
public DataPath GraphSnapshotPath = new();

[Button(ButtonSizes.Small, Name = "Load Data")]         // editor only
[ButtonGroup(DataButtonsGroup), RuntimeButton(DataButtonsGroup, "Load Data")]
public virtual async UniTask LoadDataAsync() { ... }
  \end{lstlisting}
  \vspace{1ex}

  \begin{itemize}
    \item The annotation is inert either way; only the reader differs. Odin's
      runs in the editor, \texttt{RuntimeTabMenu} in the running game, and
      neither knows about the other.
    \item Hence the pairs: \texttt{[TabGroup]}/\texttt{[RuntimeTab]},
      \texttt{[Button]}/\texttt{[RuntimeButton]},
      \texttt{[ShowIf]}/\texttt{[RuntimeShowIf]},
      \texttt{[PropertyOrder]}/\texttt{[RuntimeGroupOrder]}.
  \end{itemize}
\end{frame}

\begin{frame}[fragile]{How the Menu Reads Them}
  \begin{lstlisting}[basicstyle=\scriptsize\ttfamily]
// Assets/SEE/UI/RuntimeConfigMenu/RuntimeTabMenu.cs
city.GetType().GetMembers().Where(IsCityAttribute)     // public, of a city
    .OrderBy(HasTabAttribute).ThenBy(GetTabName)       // and two more keys
    .ForEach(memberInfo => CreateSetting(memberInfo, null, city));

// [RuntimeShowIf(nameof(Kind), ...)] keeps its condition as a string
MethodInfo m = FindMemberRecursive(type, condition, ...);  // up the chain
if (m != null && m.ReturnType == typeof(bool))   // the check the compiler
{                                                // cannot make for you
    return (bool)m.Invoke(obj, null);
}
  \end{lstlisting}
  \vspace{1ex}

  \begin{itemize}
    \item Forgetting a twin is caught only once the menu is built:
      \texttt{AssertConditionalAttributePairsValid} throws if a
      \texttt{[ShowIf]} has no \texttt{[RuntimeShowIf]} --- asked with
      \texttt{inherit: false}, so only the member's own annotation counts.
    \item \texttt{FindMemberRecursive} climbs \texttt{BaseType} itself, because
      \texttt{BindingFlags.NonPublic} finds the private members of one type,
      never the inherited ones.
  \end{itemize}
\end{frame}

\begin{frame}{What Reflection Costs}
  It buys you decisions that the compiler cannot make, and it pays for them
  with everything the compiler would have given you.
  \vspace{1ex}

  \begin{itemize}
    \item \textbf{No compile-time check.} Names are strings. Rename a method
      and the annotation still looks right, while the lookup quietly returns
      \texttt{null}. Use \texttt{nameof} wherever a name is not read from data.
    \item \textbf{A price per lookup.} Scanning an assembly and building
      \texttt{Info} objects costs orders of magnitude more than a direct call.
      Do it once and remember the result, as
      \texttt{ParsingConfigFactory} does in its static constructor. Never per
      frame.
    \item \textbf{Invisible to the linker.} A type that nothing references
      statically may be stripped from an IL2CPP build --- SEE's Android build
      is one --- and then exists in the editor but not in the player.
  \end{itemize}
\end{frame}

\begin{exerciseframe}{Fields from the Environment}
  SEE can set a field from an environment variable, so that a machine can be
  configured without touching a scene:
  \vspace{1ex}

  \begin{lstlisting}[basicstyle=\scriptsize\ttfamily]
// Assets/SEE/Utils/EnvironmentVariable.cs
[AttributeUsage(AttributeTargets.Field)]
public class EnvironmentVariableAttribute : Attribute
{
  public string VariableName { get; }
  public EnvironmentVariableAttribute(string variableName) { ... }
}

[EnvironmentVariable("DEBUG_LOG")]        // somewhere in a class
public bool debugLog;
  \end{lstlisting}
  \vspace{1ex}

  Write \texttt{SetEnvironmentVariableFields<T>(T target)}: every annotated
  field of \texttt{target} gets the value of its environment variable, and a
  field whose variable is unset stays as it is.
\end{exerciseframe}

\begin{solutionframe}{Fields from the Environment}
  \begin{lstlisting}[basicstyle=\scriptsize\ttfamily]
public static void SetEnvironmentVariableFields<T>(T target)
{
  foreach (FieldInfo field in typeof(T).GetFields())   // public ones only
  {
    string name = field.GetCustomAttribute<EnvironmentVariableAttribute>()
                       ?.VariableName;
    if (name != null && Environment.GetEnvironmentVariable(name) is { } value)
    {
      field.SetValue(target, Convert.ChangeType(value, field.FieldType));
    }
  }
}
  \end{lstlisting}
  \vspace{0.5ex}

  \begin{itemize}
    \item \texttt{GetFields()} returns \emph{public instance} fields only. A
      private one needs \texttt{BindingFlags.NonPublic |
      BindingFlags.Instance}; forget that and the annotation is ignored in
      silence.
    \item The environment hands you a \texttt{string}, so
      \texttt{Convert.ChangeType} does the rest. SEE catches
      \texttt{FormatException} and \texttt{OverflowException} around it: a typo
      in the variable is the user's mistake, not a reason to crash.
  \end{itemize}
\end{solutionframe}

\begin{exerciseframe}{A Button from an Annotation}
  SEE's runtime configuration menu shows a button for annotated methods:
  \vspace{1ex}

  \begin{lstlisting}[basicstyle=\scriptsize\ttfamily]
// Assets/SEE/UI/RuntimeConfigMenu/RuntimeButtonAttribute.cs
public class RuntimeButtonAttribute : Attribute
{
  public readonly string Name;
  public readonly string Label;
  ...
}

// Assets/SEE/Game/City/SEECityEvolution.cs
[RuntimeButton(DataButtonsGroup, "Load and Draw")]
public async UniTask StartEvolutionAsync() { ... }
  \end{lstlisting}
  \vspace{1ex}

  Write what the menu needs: collect the annotated methods of a city and call
  the one whose button was clicked. Which methods must you refuse, and when do
  you learn that you should have?
\end{exerciseframe}

\begin{solutionframe}{A Button from an Annotation}
  \begin{lstlisting}[basicstyle=\scriptsize\ttfamily]
foreach (MethodInfo method in city.GetType().GetMethods())
{
  RuntimeButtonAttribute button
      = method.GetCustomAttribute<RuntimeButtonAttribute>();
  if (button != null && method.GetParameters().Length == 0)  // refuse the
  {                                                          // others
    AddButton(button.Label, () => method.Invoke(city, null));
  }
}
  \end{lstlisting}
  \vspace{0.5ex}

  \begin{itemize}
    \item An attribute is inert data: it does nothing until somebody looks for
      it. In SEE that somebody is \texttt{RuntimeTabMenu}, and it also refuses
      methods that expect arguments.
    \item \texttt{Invoke} passes its arguments as \texttt{object[]}, so nothing
      here is checked at compile time. A wrong signature, a renamed method, a
      typo in the label --- all of them surface when the menu is opened, not
      when the code is built.
  \end{itemize}
\end{solutionframe}

\begin{exerciseframe}{A Factory Without a \texttt{switch}}
  Every metric tool that SEE reads brings its own configuration class:
  \vspace{1ex}

  \begin{lstlisting}[basicstyle=\scriptsize\ttfamily]
// Assets/SEE/DataModel/DG/IO/
public abstract class ParsingConfig
{
  public string ToolId = string.Empty;    // "JaCoCo", "Checkstyle", ...
}

public class JaCoCoParsingConfig     : ParsingConfig { ... }
public class CheckstyleParsingConfig : ParsingConfig { ... }
public class MSBuildParsingConfig    : ParsingConfig { ... }
  \end{lstlisting}
  \vspace{1ex}

  \texttt{ParsingConfigFactory.Create("JaCoCo")} shall return the matching
  configuration, and adding a tool tomorrow shall need no change to the
  factory. Write the registration.
\end{exerciseframe}

\begin{solutionframe}{A Factory Without a \texttt{switch}}
  \begin{lstlisting}[basicstyle=\scriptsize\ttfamily]
static ParsingConfigFactory()            // runs once, on first access
{
  IEnumerable<Type> types = Assembly.GetExecutingAssembly().GetTypes()
      .Where(t => t.IsSubclassOf(typeof(ParsingConfig)) && !t.IsAbstract);
  foreach (Type type in types)
  {
    ParsingConfig probe = (ParsingConfig)Activator.CreateInstance(type);
    registry.Add(probe.ToolId, type);    // SEE reports duplicate ToolIds
  }
}
  \end{lstlisting}
  \vspace{0.5ex}

  \begin{itemize}
    \item \textbf{The instance comes first, the key second.} \texttt{ToolId} is
      an instance field, so the factory builds a probe to read it, and
      \texttt{Activator.CreateInstance} wants a public parameterless
      constructor. A class without one throws while the factory scans, where
      SEE catches it per type and logs it rather than losing the registry.
    \item \textbf{The price is paid once}, in the static constructor. Afterwards
      \texttt{Create} is a dictionary lookup, and a new tool needs nothing but
      a new subclass.
  \end{itemize}
\end{solutionframe}
